\chapter{System Design and Implementation}
\label{ch:system_design}

\section{End-to-End System Architecture}
\label{sec:design_architecture}

\subsection{High-Level Architecture Overview}
\label{subsec:design_arc_overview}

The framework follows a hub-and-spoke topology in which a central coordinator
orchestrates federated learning, root cause analysis, and model versioning,
while distributed edge agents perform local telemetry collection, inference,
and model training. Figure~\ref{fig:architecture_overview} presents the
high-level architecture.

\begin{figure}[H]
\centering
\begin{tikzpicture}[
    node distance=1.2cm and 2.0cm,
    box/.style={draw, rounded corners, minimum width=2.8cm, minimum height=0.9cm,
                align=center, font=\small},
    edge_box/.style={box, fill=blue!8},
    coord_box/.style={box, fill=orange!10},
    data_box/.style={box, fill=green!8},
    arrow/.style={-{Stealth[length=2.5mm]}, thick},
]
% Edge Nodes
\node[edge_box] (e1) {Edge Agent 1\\[-2pt]\scriptsize LSTM-AE + Q-Thresh};
\node[edge_box, below=0.6cm of e1] (e2) {Edge Agent 2\\[-2pt]\scriptsize LSTM-AE + Q-Thresh};
\node[edge_box, below=0.6cm of e2] (e3) {Edge Agent $K$\\[-2pt]\scriptsize LSTM-AE + Q-Thresh};

% Coordinator
\node[coord_box, right=3.5cm of e2] (coord) {Central\\Coordinator};
\node[coord_box, above=0.7cm of coord] (fl) {FL Aggregator\\[-2pt]\scriptsize FedAvg + DP};
\node[coord_box, below=0.7cm of coord] (rca) {RCA Engine\\[-2pt]\scriptsize PageRank};

% Data stores
\node[data_box, right=2.5cm of coord] (db) {Model Registry\\[-2pt]\scriptsize SQLite / WAL};
\node[data_box, right=2.5cm of fl] (graph) {Causal Graph\\[-2pt]\scriptsize NetworkX};

% Traces
\node[data_box, below=0.7cm of rca] (traces) {Trace Collector\\[-2pt]\scriptsize OpenTelemetry};

% Edges -> Coordinator
\draw[arrow] (e1.east) -- node[above, font=\scriptsize, sloped] {gRPC updates} (fl.west);
\draw[arrow] (e2.east) -- (coord.west);
\draw[arrow] (e3.east) -- node[below, font=\scriptsize, sloped] {anomaly events} (rca.west);

% Coordinator internals
\draw[arrow] (fl) -- (coord);
\draw[arrow] (rca) -- (coord);
\draw[arrow] (coord) -- (db);
\draw[arrow] (fl) -- (graph);
\draw[arrow] (traces) -- (rca);

% Broadcast back
\draw[arrow, dashed] (coord.west) -- node[above, font=\scriptsize, sloped] {global model} (e2.east);
\end{tikzpicture}
\caption{High-level architecture of the distributed anomaly detection framework.
         Edge agents communicate with the central coordinator via gRPC.
         Arrows indicate primary data flows; dashed lines indicate broadcast.}
\label{fig:architecture_overview}
\end{figure}

The architecture comprises five principal subsystems:

\begin{enumerate}
    \item \textbf{Edge Agents}: Deployed alongside each monitored microservice.
    Each agent collects local telemetry, runs inference using a local LSTM-AE
    model, applies the adaptive Q-learning threshold, and participates in
    federated training rounds.

    \item \textbf{Central Coordinator}: Orchestrates FL rounds, hosts the RCA
    engine, and manages the global model registry. It is designed as a
    stateless process backed by durable storage, enabling active-passive
    failover (Chapter~\ref{ch:reliability}).

    \item \textbf{FL Aggregator}: Receives compressed model updates from edge
    agents, validates them, and computes the FedAvg-aggregated global model.
    Differential privacy noise is injected at the client side before
    transmission (Chapter~\ref{ch:federated_learning}).

    \item \textbf{RCA Engine}: Maintains the causal dependency graph and, upon
    receiving batched anomaly events, ranks candidate root causes using the
    scoring approach described in Chapter~\ref{ch:root_cause_analysis}.

    \item \textbf{Trace Collector}: Ingests distributed traces from
    OpenTelemetry exporters and continuously updates the causal graph with
    inter-service dependency observations.
\end{enumerate}

\subsection{How Components Fit Together}
\label{subsec:design_arc_integration}

A complete detection cycle proceeds through the following stages:

\begin{enumerate}
    \item \textbf{Metric Collection.} The edge agent scrapes local metrics
    (CPU, memory, latency, error rate) at a configurable interval (default
    10\,s) and writes them into a sliding-window buffer of length $W=50$.

    \item \textbf{Inference.} The LSTM-AE processes the current window and
    outputs a reconstruction error $e_t^{(i)}$.

    \item \textbf{Threshold Decision.} The per-service Q-learning threshold
    $\tau_t^{(i)}$ is compared against $e_t^{(i)}$. If exceeded, the agent
    emits an anomaly event to the coordinator.

    \item \textbf{RCA Query.} The coordinator batches anomaly events within a
    60\,s window and invokes the RCA engine. The engine classifies anomalous
    services as root causes or propagated anomalies and returns a ranked list.

    \item \textbf{Alert Enrichment.} The RCA ranking, propagation paths, and
    graph metadata are attached to the alert and forwarded to the alerting
    layer (PagerDuty, Slack, or webhook).

    \item \textbf{Federated Learning.} In a separate background loop, the
    coordinator schedules FL rounds. Edge agents train locally, apply gradient
    clipping and optional DP noise, compress the update, and transmit it via
    gRPC. The coordinator aggregates updates using FedAvg and broadcasts the
    new global model.

    \item \textbf{Threshold Adaptation.} After each detection decision, the
    outcome (TP/FP/FN/TN) is fed back to the Q-learning agent, which
    periodically adjusts the threshold via the Bellman update
    (Chapter~\ref{ch:adaptive_thresholding}).
\end{enumerate}

\section{Data Flow}
\label{sec:design_dataflow}

\subsection{Metric Ingestion Pipeline}
\label{subsec:design_ingestion}

Each edge agent runs a lightweight metric collector that interfaces with
Prometheus exporters or StatsD receivers deployed alongside the target
microservice. Raw metrics are buffered in a fixed-capacity ring buffer (default
1\,000 observations) and preprocessed as follows:

\begin{enumerate}
    \item \textbf{Normalization.} Each feature is z-score normalized using a
    running mean and standard deviation computed over the last $N=500$
    observations.

    \item \textbf{Windowing.} The normalized time series is segmented into
    overlapping sliding windows of length $W=50$ with a stride of~1.

    \item \textbf{Tensor Construction.} Each window is reshaped into a
    three-dimensional PyTorch tensor of shape
    $(\text{batch}=1, W, \text{features})$ for input to the LSTM-AE.
\end{enumerate}

\subsection{Anomaly Detection Pipeline}
\label{subsec:design_detection_pipeline}

The anomaly detection pipeline operates per-window as follows:

\begin{enumerate}
    \item The LSTM-AE encoder maps the input window $\mathbf{x} \in
    \mathbb{R}^{W \times F}$ to a latent representation via a two-layer LSTM
    with hidden size~64.

    \item The decoder reconstructs $\hat{\mathbf{x}}$ from the latent state.

    \item The reconstruction error is computed as the mean absolute error
    reduced over the sequence and feature dimensions:
    $e = \frac{1}{WF}\sum_{t,f}|x_{t,f} - \hat{x}_{t,f}|$.

    \item The error is compared against the adaptive threshold
    $\tau_t^{(i)}$. If $e > \tau_t^{(i)}$, an anomaly event is emitted.
\end{enumerate}

Inference latency is approximately 0.07\,ms per window on CPU
(Section~\ref{sec:eval_detection}), ensuring that real-time detection does not
bottleneck the metric collection loop.

\subsection{Federated Learning Pipeline}
\label{subsec:design_fl_pipeline}

Each FL round follows the pipeline illustrated below.

\begin{enumerate}
    \item \textbf{Round Scheduling.} The coordinator initiates a round every
    $T_{\text{round}}$ seconds (default 300\,s) or after receiving a
    configurable number of new local training samples.

    \item \textbf{Local Training.} Each participating edge agent trains for
    $E=5$ local epochs using the Adam optimiser ($\text{lr}=10^{-3}$,
    batch size~64) on its local dataset.

    \item \textbf{DP Noise Injection.} If differential privacy is enabled,
    per-sample gradients are clipped to norm $C=1.0$ and Gaussian noise
    $\mathcal{N}(0, \sigma^2 C^2\mathbf{I})$ is added, where
    $\sigma$ is the calibrated noise multiplier corresponding to the
    desired privacy budget $\varepsilon$.

    \item \textbf{Serialization.} The model state dictionary is serialized
    using \texttt{torch.save()} (binary) or Protocol Buffers for
    cross-language interoperability.

    \item \textbf{Compression.} The serialized payload is compressed using the
    adaptive compression controller: Zstd is selected when CPU utilization is
    below 80\%, and LZ4 is selected otherwise
    (Section~\ref{subsec:eval_adaptive_compression}).

    \item \textbf{Transmission.} The compressed payload is uploaded to the
    coordinator via the \texttt{SubmitUpdate} gRPC endpoint.

    \item \textbf{Aggregation.} The coordinator computes the FedAvg-weighted
    average of received updates and persists the new global model.

    \item \textbf{Broadcast.} Edge agents poll for the updated model via
    \texttt{GetGlobalModel} and atomically swap their local weights.
\end{enumerate}

\subsection{Root Cause Analysis Pipeline}
\label{subsec:design_rca_pipeline}

The RCA pipeline connects trace ingestion to alert enrichment:

\begin{enumerate}
    \item \textbf{Trace Ingestion.} The Trace Collector receives spans from
    OpenTelemetry or Jaeger exporters and extracts parent--child
    inter-service call records.

    \item \textbf{Dependency Graph Update.} Each call record updates the
    corresponding edge weight in the {\texttt{CausalGraph}} using exponential
    decay (Section~\ref{subsec:rca_dependency_graph}).

    \item \textbf{Anomaly Event Seed.} When anomaly events arrive, the
    affected services are marked as anomalous in the graph.

    \item \textbf{Root Cause Scoring.} The \texttt{RootCauseAnalyzer}
    classifies anomalous services as root causes or propagated anomalies
    based on upstream dependency analysis and impact scoring.

    \item \textbf{Alert Enrichment.} The top-$N$ ranked candidates, their
    scores, and explanatory propagation paths are attached to the anomaly
    alert.
\end{enumerate}

\section{Component Design}
\label{sec:design_components}

\subsection{Edge Node Agent}
\label{subsec:design_edge_agent}

The edge agent is a single-process Python application with three concurrent
execution loops:

\begin{itemize}
    \item \textbf{Inference Loop} (main thread): Continuously reads from the
    metric buffer, runs LSTM-AE inference, applies the Q-learning threshold,
    and emits anomaly events.

    \item \textbf{FL Training Loop} (background thread): Triggered by FL round
    notifications. Downloads the global model, trains locally, applies DP, and
    uploads the compressed update.

    \item \textbf{Threshold Adaptation Loop}: Every 10~observations, invokes
    \texttt{ThresholdTuner.tune\_threshold()} to adjust the service-specific
    threshold.
\end{itemize}

The agent is configured via environment variables and a \texttt{config.yaml}
file, supporting two profiles: \emph{lightweight} (stateless) and
\emph{standard} (persistent SQLite cache), as described in
Chapter~\ref{ch:federated_learning}.

\subsection{Central Coordinator}
\label{subsec:design_coordinator}

The coordinator exposes a gRPC service with the following endpoints:

\begin{itemize}
    \item \texttt{RegisterNode}: Registers an edge agent with its node
    profile, hardware capabilities, and initial sample count.

    \item \texttt{SubmitUpdate}: Receives a compressed model update. The
    coordinator decompresses, validates (optional poisoning detection), and
    buffers the update.

    \item \texttt{GetGlobalModel}: Returns the latest global model checkpoint
    along with the current round number and compression metadata.

    \item \texttt{ReportAnomaly}: Receives anomaly events from edge agents
    for RCA processing.
\end{itemize}

Persistent state---including the model registry (last $N=10$ checkpoints), the
causal graph snapshots, and the FL round log---is stored in an SQLite database
with Write-Ahead Logging (WAL) enabled for crash-safe durability.

\subsection{Trace Collector}
\label{subsec:design_trace_collector}

The Trace Collector is an asynchronous module that:

\begin{enumerate}
    \item Subscribes to OpenTelemetry OTLP/gRPC or Jaeger Thrift trace
    streams.

    \item Parses spans into inter-service dependency records.

    \item Feeds records to the \texttt{CausalGraph.update\_from\_traces()}
    method, which updates edge weights, applies decay, and prunes weak edges.

    \item Triggers periodic graph snapshots (every 300\,s) for persistence.
\end{enumerate}

\subsection{Alerting and Dashboarding}
\label{subsec:design_alerting}

Alerts are structured JSON payloads containing:

\begin{itemize}
    \item \texttt{service}: Affected service name.
    \item \texttt{anomaly\_score}: Reconstruction error value.
    \item \texttt{threshold}: Current adaptive threshold.
    \item \texttt{rca\_candidates}: Ranked list of root-cause services with
    impact scores.
    \item \texttt{propagation\_paths}: Explanatory dependency paths.
    \item \texttt{timestamp}: Detection timestamp.
\end{itemize}

Alerts are delivered via configurable sinks: PagerDuty, Slack webhooks, or
generic HTTP endpoints. A Grafana dashboard provides real-time visualization of
anomaly scores, FL round status, causal graph topology, and threshold history
per service.

\section{System Orchestration}
\label{sec:design_orchestration}

\subsection{Kubernetes Deployment}
\label{subsec:design_k8s}

The framework is designed for deployment on Kubernetes, with the following
resource topology:

\begin{itemize}
    \item \textbf{Coordinator}: Deployed as a \texttt{Deployment} with a
    single replica (active-passive failover via leader election). A
    \texttt{PersistentVolumeClaim} provides durable storage for the SQLite
    database. Horizontal Pod Autoscaling (HPA) is configured to scale based
    on gRPC request rate if the coordinator is deployed as a stateless
    service backed by an external database.

    \item \textbf{Edge Agents}: Deployed as a \texttt{DaemonSet} to ensure one
    agent per monitored node, or as a \texttt{Deployment} with pod anti-affinity
    to co-locate agents with target microservice pods.

    \item \textbf{Service Discovery}: Edge agents discover the coordinator via
    Kubernetes DNS \\ (\texttt{coordinator-svc.namespace.svc.cluster.local}).

    \item \textbf{Configuration}: FL hyperparameters, privacy budget, and
    compression settings are supplied via \texttt{ConfigMaps}. Secrets (e.g.,
    gRPC TLS certificates) are stored in Kubernetes \texttt{Secrets}.
\end{itemize}

\subsection{Configuration Management}
\label{subsec:design_config}

All runtime parameters are managed through a hierarchical
\texttt{config.yaml} file, structured as follows:

\begin{verbatim}
model:
  input_size: 10
  hidden_size: 64
  num_layers: 2
  dropout: 0.2
federated:
  rounds: 100
  local_epochs: 5
  min_clients: 3
  round_timeout_s: 120
privacy:
  enabled: true
  epsilon: 1.0
  delta: 1e-5
  clip_norm: 1.0
compression:
  default_algorithm: zstd
  cpu_threshold: 80
thresholding:
  learning_rate: 0.1
  discount_factor: 0.95
  epsilon: 0.1
\end{verbatim}

The configuration module supports live reload: the coordinator watches the
config file for changes and applies non-destructive updates (e.g., adjusting
$\varepsilon$ or compression level) without restarting the process.

\subsection{Managing Data Flow and Scheduling}
\label{subsec:design_scheduling}

FL rounds are triggered by one of two policies:

\begin{itemize}
    \item \textbf{Time-based}: The coordinator initiates a new round every
    $T_{\text{round}}$ seconds (default 300\,s). This is suitable for
    steady-state operation.

    \item \textbf{Event-based}: A round is triggered after a configurable
    number of new local training samples ($M=1000$) have been collected across
    registered nodes. This adaptive policy avoids unnecessary rounds during
    low-traffic periods.
\end{itemize}

The coordinator maintains a round log recording the start time, participating
nodes, aggregation result, and round duration. Overlapping rounds are
prevented by a round-state lock: if a round is in progress, new triggers are
deferred until the current round completes or times out.

\section{Implementation Details}
\label{sec:design_implementation}

\subsection{Technology Stack}
\label{subsec:design_stack}

Table~\ref{tab:tech_stack} summarizes the technology stack.

\begin{table}[H]
    \centering
    \begin{tabular}{|p{0.28\textwidth}|p{0.62\textwidth}|}
        \hline
        \textbf{Component} & \textbf{Technology} \\
        \hline
        Language & Python 3.12 \\
        \hline
        ML Framework & PyTorch 2.x (LSTM-AE, DP noise injection) \\
        \hline
        RPC Framework & gRPC (grpcio + protobuf) \\
        \hline
        Compression & python-lz4, zstandard \\
        \hline
        Graph Library & NetworkX (causal graph, BFS, impact scoring) \\
        \hline
        Persistent Storage & SQLite with WAL mode (via Python \texttt{sqlite3}) \\
        \hline
        Caching / Buffering & Redis (optional, for update queuing) \\
        \hline
        Monitoring & Prometheus client, OpenTelemetry SDK \\
        \hline
        Container Runtime & Docker \\
        \hline
        Orchestration & Kubernetes (Minikube for local development) \\
        \hline
        Chaos Engineering & Chaos Mesh (fault injection for benchmarking) \\
        \hline
        Load Testing & Locust (traffic generation) \\
        \hline
    \end{tabular}
    \caption{Technology stack used in the framework implementation}
    \label{tab:tech_stack}
\end{table}

\subsection{Key Implementation Decisions}
\label{subsec:design_decisions}

Several design decisions were made with explicit rationale:

\begin{itemize}
    \item \textbf{SQLite over in-memory stores.} SQLite with WAL mode provides
    ACID-compliant persistence with zero-configuration deployment. Unlike Redis
    or PostgreSQL, it requires no external process, making it suitable for
    single-coordinator deployments. WAL mode ensures that readers and writers
    do not block each other, and that in-progress writes survive crashes
    without corrupting the database \cite{ongaro2014raft}.

    \item \textbf{PyTorch over TensorFlow.} PyTorch provides imperative
    execution, which simplifies debugging of the LSTM-AE training loop and
    gradient manipulation (clipping, noise injection) required for differential
    privacy. Additionally, PyTorch supports Metal Performance Shaders (MPS)
    for Apple Silicon acceleration, enabling development and testing on macOS
    hardware.

    \item \textbf{gRPC over REST.} gRPC provides strongly-typed service
    contracts via Protocol Buffers, binary-encoded payloads, HTTP/2
    multiplexing, and bidirectional streaming---all of which reduce latency
    and connection overhead compared to JSON-over-HTTP REST endpoints
    \cite{grpc2016whitepaper}.

    \item \textbf{Zstd as default compression with LZ4 fallback.} Benchmark
    results (Section~\ref{subsec:eval_compression}) demonstrate that Zstd
    achieves a 10\% payload reduction on binary model weights and up to 61\%
    on JSON payloads, at the cost of higher CPU usage. The adaptive controller
    automatically switches to LZ4 when CPU utilization exceeds 80\%, preserving
    compression benefits without impacting inference performance on
    resource-constrained nodes.
\end{itemize}
